% ====================================================================
% conclusion.tex (dissertation example)
% ====================================================================

\section{Conclusion}
\label{sec:conclusion}

The aim of this project was to establish how much input noise an ideal 12-bit
converter needs for oversampling and averaging to deliver the resolution that
theory predicts. Taking the objectives in turn:

\begin{enumerate}
  \item The theory was reviewed and reduced to a single testable prediction,
        \cref{eq:enobTheory}, which gives the effective resolution in terms
        of the oversampling ratio and the input noise.
  \item A simulation model was built and checked against the case with no
        averaging, where it gave 10.97 bits against a predicted 11.00.
  \item The effective resolution was measured for oversampling ratios from 1
        to 256, for a sine wave and for static inputs at four noise levels.
  \item With \qty{0.5}{LSB} rms of noise or more, the results agreed with the
        theory for both kinds of input, reaching 15.0 effective bits at a
        ratio of 256. With \qty{0.25}{LSB} the improvement stopped at around
        14 bits, and with no noise there was none for a static input.
\end{enumerate}

The answer to the question is therefore that about half an LSB rms of noise
is sufficient up to an oversampling ratio of 256, and that this is also close
to the best choice, as doubling it cost almost a bit of resolution. This is a
more specific figure than the one LSB or more given in the guidance that was
reviewed.

\subsection{Further work}

The most useful next step would be to repeat the static test on a real
microcontroller, using a precision voltage source for the input, to find how
far the non-linearity of an actual converter limits what the simulation
predicts.

The model could also be extended in two directions. The first is to replace
the Gaussian noise with a deliberately added triangular waveform, which is
simpler to generate in a circuit and may need a smaller amplitude for the
same result. The second is to carry the oversampling ratio beyond 256, as the
curve for \qty{0.5}{LSB} must eventually level out in the same way as the
curve for \qty{0.25}{LSB}, and it would be useful to know where.
